\par\Needspace{6\baselineskip}\medskip
\begin{center}\hypertarget{native-243-article173}{}\pdfbookmark[2]{173.}{article-173}173.\end{center}
Maximi imprimis momenti est, ut singulae corporis coelestis positiones
geocentricae, quibus orbitam superstruere propositum est, non ex observationibus
solitariis petitae sint, sed si fieri potest e pluribus ita combinatis, ut
errores forte commissi quantum licet sese mutuo destruxerint. Observationes
scilicet tales, quae paucorum dierum intervallo ab invicem distant --- vel adeo
prout res fert intervallo 15 aut 20 dierum --- in calculo non adhibendae erunt
tamquam totidem positiones diversae, sed potius positio unica inde derivabitur,
quae inter cunctas quasi media est, adeoque praecisionem longe maiorem
admittit, quam observationes singulae seorsim consideratae. Quod negotium
sequentibus principiis innititur.

Corporis coelestis loca geocentrica ex elementis approximatis calculata a
locis veris parum discrepare, differentiaeque inter haec et illa mutationes
lentissimas tantum subire debent, ita ut intra paucorum dierum decursum
propemodum pro constantibus haberi queant, vel saltem variationes tamquam
temporibus proportionales spectandae sint. Si itaque observationes ab omni
errore immunes essent, differentiae inter locos observatos temporibus $t$, $t'$, $t''$, $t'''$
etc. respondentes eosque qui ex elementis computati sunt, i.\,e. differentiae
tum longitudinum tum latitudinum, sive tum ascensionum rectarum tum
declinationum, observatarum a computatis, forent quantitates vel sensibiliter
aequales, vel saltem uniformiter lentissimeque increscentes aut decrescentes.
Respondeant e.\,g. illis temporibus ascensiones rectae observatae $a$, $a'$, $a''$, $a'''$
etc., computatae autem sint $a+\delta$, $a'+\delta'$, $a''+\delta''$, $a'''+\delta'''$
etc.; tunc differentiae $\delta$, $\delta'$, $\delta''$, $\delta'''$ etc. a veris
elementorum deviationibus eatenus tantum discrepabunt, quatenus observationes
ipsae sunt erroneae: si itaque illas deviationes pro omnibus istis
observationibus tamquam constantes spectare licet, exhibebunt quantitates
$\delta$, $\delta'$, $\delta''$, $\delta'''$ etc. totidem determinationes diversas
eiusdem magnitudinis, pro cuius valore correcto itaque assumere conveniet
medium arithmeticum inter illas determinationes, quatenus quidem nulla adest
ratio, cur unam \SourceLeaf{244}alteramve praeferamus. Sin vero observationibus
singulis idem praecisionis gradus haud attribuendus videtur, supponamus
praecisionis gradum in singulis resp. proportionalem aestimandum esse numeris
$e$, $e'$, $e''$, $e'''$ etc., i.\,e. errores his numeris reciproce proportionales
in observationibus aeque facile committi potuisse; tum secundum principia
infra tradenda valor medius maxime probabilis haud amplius erit medium
arithmeticum simplex, sed
\[
=\frac{ee\delta+e'e'\delta'+e''e''\delta''+e'''e'''\delta'''+\text{etc.}}
{ee+e'e'+e''e''+e'''e'''+\text{etc.}}\mathpunct{\cdot}
\]
Statuendo iam hunc valorem medium $=\Delta$, pro ascensionibus rectis veris
assumere licebit resp. $a+\delta-\Delta$, $a'+\delta'-\Delta$, $a''+\delta''-\Delta$,
$a'''+\delta'''-\Delta$ etc., tumque arbitrarium erit, quanam in calculo utamur.
Quodsi vero vel observationes temporis intervallo nimis magno ab invicem
distant, aut si orbitae elementa satis approximata nondum innotuerant, ita ut
non licuerit, horum deviationes tamquam constantes pro observationibus cunctis
spectare, facile perspicietur, aliam hinc differentiam non oriri, nisi quod
deviatio media sic inventa non tam omnibus observationibus communis
supponenda erit, quam potius ad tempus medium quoddam referenda, quod perinde
e singulis temporum momentis derivare oportet, ut $\Delta$ ex singulis
deviationibus, adeoque generaliter ad tempus
$\frac{eet+e'e't'+e''e''t''+e'''e'''t'''+\text{etc.}}{ee+e'e'+e''e''+e'''e'''+\text{etc.}}\mathpunct{\cdot}$
Si itaque summam praecisionem appetere placet, pro eodem tempore locum
geocentricum ex elementis computare, ac dein ab errore medio $\Delta$ liberare
oportebit, ut positio quam accuratissima emergat: plerumque tamen abunde
sufficiet, si error medius ad observationem tempori medio proximam referatur.
Quae hic de ascensionibus rectis diximus, perinde de declinationibus, aut si
mavis de longitudinibus et latitudinibus valent: attamen semper praestabit,
immediate ascensiones rectas et declinationes ex elementis computatas cum
observatis comparare; sic enim non modo calculum magis expeditum lucramur,
praesertim si methodis in artt.\ 53---60 expositis utimur, sed eo insuper
titulo illa ratio se commendat, quod observationes incompletas quoque in usum
vocare licet, praetereaque, si omnia ad longitudines et latitudines
referrentur, metuendum esset, ne observatio quoad ascensionem recte, quoad
declinationem male instituta (vel vice versa) ab utraque parte depravetur,
atque sic prorsus inutilis evadat. --- Ceterum gradus praecisionis medio ita
invento attribuendus secundum principia mox explicanda erit
$=\surd(ee+e'e'+e''e''+e'''e'''+\text{etc.})$, ita ut quatuor vel novem
observationes aeque \SourceLeaf{245}exactae requirantur, si medium praecisione
dupla vel tripla gaudere debet, et sic porro.
